Multi-agent systems are increasingly used for forecasting future events, on the assumption that deliberation among multiple LLMs improves reasoning and calibration.
Existing approaches, however, overlook a basic design choice: \emph{what information each agent receives}.
When all agents observe identical evidence, deliberation collapses into herding rather than genuine belief revision, and the ensemble offers little over a single agent.
We propose \emph{designed information asymmetry}: partitioning the evidence into a shared public pool and disjoint private subsets, so that each agent holds exclusive knowledge that can reach the others only through deliberation.
We show theoretically that this partition reduces inter-agent error correlation, and instantiate it in \textbf{InfoDelphi}, a framework combining relevance-aware evidence routing, rationale-based iterative deliberation, and a calibrated aggregation layer that counteracts systematic LLM overconfidence via asymmetric shrinkage.
On \textsc{PolyGym}, a controlled benchmark of 250 resolved Polymarket questions with fixed pre-retrieved evidence, and on 231 questions from FutureX, InfoDelphi attains Brier scores of 0.158 and 0.201, outperforming nine single-agent and multi-agent baselines under a paired bootstrap test on both datasets, with a 13--18\% relative Brier improvement over the strongest baseline.
Analysis shows that deliberation under homogeneous input collapses forecast diversity while partitioned evidence preserves it, and that overconfidence is best corrected in the aggregation layer: prompt-level skepticism depresses the entire forecast distribution and hurts accuracy.
